\beginsong{La Blanche Hermine}[by={Gilles Servat}]
\beginverse
\[Am]J'ai rencontré ce matin
Devant la \[G]haie de mon champ
Une \[F]troupe de marins
D'ou\[Em]vriers, de pay\[Am]sans.
\endverse
\beginverse
\[Am]Où allez-vous camarades
Avec \[G]vos fusils chargés ?
Nous \[F]tendrons des embuscades,
Viens \[Em]rejoindre notre ar\[Am]mée.
\endverse
\beginchorus
\[Am]La voilà la Blanche Hermine,
Vive la \[G]mouette et l'ajonc !
La voi\[F]là la Blanche Hermine,
Vive \[Em]Fougères et Clis\[Am]son !
\endchorus
\beginverse
\[Am]Ma mie dit que c'est folie
D'aller \[G]faire la guerre aux Francs.
Moi \[F]je dis que c'est folie
\[Em]D'être enchaîné plus long\[Am]temps.
\endverse
\beginverse
\[Am]Elle aura bien de la peine
Pour \[G]élever les enfants.
Elle \[F]aura bien de la peine
Car \[Em]je m'en vais pour long\[Am]temps.
\endverse
\beginverse
\[Am]Et si je meurs à la guerre,
Pourra-\[G]t-elle me pardonner
D'avoir \[F]préféré ma terre
À l'a\[Em]mour qu'elle me don\[Am]nait ?
\endverse
\endsong
